% Explain how we do error handling.

An essential part of making a programming language usable includes having a sensible error handling.
Clear error messages help users understand what went wrong in their program and how to fix it.
To facilitate this, we perform error handling in all parts of the interpreter.

At the center of the error handling is the \texttt{VVPLController} class which functions as a centralized hub for
reporting errors. This design decouples the error \emph{detection} from the error \emph{reporting}.
Each component such as the scanner or parser is responsible for identifying errors and notifying the
\texttt{VVPLController} class which in turn is responsible for handling how these errors are stored
and reported to the user.

This interplay between the classes is facilitated by the \texttt{error} method in the controller.
\begin{lstlisting}[language=java]
public static void error(int line, String errorType, String message) {
    hadError = true;
    errorMessages.add(new ErrorMessage(line, errorType, message));
}
\end{lstlisting}

This method takes in the line where the error occured, the type of error and the error message we are reporting to the user. 
The error type is one of the different errors from the \texttt{ErrorTypeStrings} class, which allows us to categorize the errors.
The method firstly sets the \texttt{hadError} flag, which signals that the current stage of compilation has failed.
After each stage of the interpreter pipeline, the controller checks this flag. If an error has been registered the pipeline stops.
Halting the interpreter as soon as a single stage fails is part of a fail fast strategy that prevents later stages from operating 
on invalid data which would cause a cascade of confusing and possibly misleading errors.
%TODO: Skal denne sætning om at scanner, parser etc. reporter alle errors. Kan måske i stedet bruges senere til naturlig overgang.
While the pipeline stops between stages, each individual stage is designed to detect as many errors as possible before terminating,
providing the user with a comprehensive report of their errors in that stage.
This gives the user a more complete overview of issues which decreases time spent debugging by
reducing the number of compiles needed to find and fix all errors in the program.

After setting the \texttt{hadError} flag the error is added to the errorMessages list which stores all the errors.
This is a list of \texttt{ErrorMessage} objects:
\begin{lstlisting}[language=java]
private static List<ErrorMessage> errorMessages = new ArrayList<>();
\end{lstlisting}

The \texttt{ErrorMessage} class stores errors as structured data rather than just plain strings.
It stores the line number, error type and message in seperate fields which allows for easy post-processing.
Specifically this allows for sorting the 
errors by line number such that the user receives a logical and actionable list of issues in the order they appear in the code. 
This is achieved by having \texttt{ErrorMessage} implement the \texttt{Comparable} interface allowing the controller to sort the list with
the \texttt{sort} method from the Collections class.
The class also overrides the \texttt{toString} method such that the error messages can easily be reported to the
user in the format: \texttt{<error-type>, line <lineNo> <message>}.

% Kan måske passe ind et sted?:
% Our system monitors the state of the source code's correctness at each stage and logs deviations in a structured user-friendly manner.

